\documentclass{article}

\usepackage[utf8]{inputenc}
\usepackage[T1]{fontenc}
\usepackage{amsmath,amssymb,amsthm}
\usepackage{mathtools}
\usepackage{hyperref}
\usepackage{booktabs}

\makeatletter
\newcommand{\citep}[2][]{%
  \begingroup
  \def\@tmpkey{#2}%
  \ifx\@tmpkey\@empty\else
    (\ifx\@empty#1\@empty\else #1, \fi
     \@nameuse{bibkey@\@tmpkey})%
  \fi
  \endgroup
}
\expandafter\def\csname bibkey@lasser2026exo\endcsname{Exogenous~Verification}
\expandafter\def\csname bibkey@lasser2026phase\endcsname{Phase~Redundancy}
\expandafter\def\csname bibkey@lasser2026micro\endcsname{Microfoundation}
\makeatother

\newtheorem{definition}{Definition}
\newtheorem{lemma}{Lemma}
\newtheorem{theorem}{Theorem}
\newtheorem{remark}{Remark}
\newtheorem{assumption}{Assumption}

\newcommand{\Aadv}{A_{\mathrm{adv}}}
\newcommand{\deltaadv}{\delta_{\mathrm{adv}}}
\newcommand{\deltastar}{\delta_{*}}
\newcommand{\rhogap}{\rho_{\mathrm{gap}}}
\newcommand{\Pproxy}{P}
\newcommand{\Ttruth}{T}
\newcommand{\epsgap}{\varepsilon_{\mathrm{gap}}}
\newcommand{\epsnonres}{\varepsilon_{\mathrm{gap}}^{\mathrm{nonres}}}
\newcommand{\epsfloor}{\varepsilon_{\mathrm{floor}}^{\mathrm{res}}}

\newcommand{\Tbeta}{T_{\beta}}
\newcommand{\Ncasc}{N_{\mathrm{cascade}}}
\newcommand{\Nev}{N_{\mathrm{events}}}
\newcommand{\bclk}{\beta_{\mathrm{clk}}}

\title{Lemma 6 ($h_{\mathrm{detect}}$ intensivity): Full Proof Draft v4}
\author{Paper 10 working draft for codex re-review}
\date{}

\begin{document}

\maketitle

\section{Goal and changes from v3}

V3 separated ``integration horizon'' (replaced by $\Tbeta$) from
``cascade lead-time'' (still uses $\tau_{\mathrm{meta}} \leq
T_{\mathrm{cascade}}$), the right structural fix.  Round~C confirmed
the structural separation but flagged five issues, the most
important being that v3's $\Tbeta = (\log(1/\alpha) +
\log(1/\beta))/(\kappa \deltastar)$ does not match Lemma~4's
actual Hoeffding form in the repo.

\textbf{Round C significant problem (Q1 algebra).}  Lemma~4
(\texttt{appendices/proofs.tex} §A.4) gives the Hoeffding form
$\Pr[T_{\mathrm{detect}} > t] \leq \exp(-2(t\delta - A)^2/(tR^2))$,
asymptotically $\exp(-2T\delta^2/R^2) =: \exp(-\kappa T \delta)$
with $\kappa = 2\delta/R^2$.  Therefore $\kappa \deltastar$ scales as
$\deltastar^2$ (not $\deltastar$), and v3's $\Tbeta$ formula is
algebraically wrong.  V4 derives the correct quantile directly from
the Hoeffding form: $\Tbeta = R^2 \log(1/\beta)/(2 \deltastar^2)$
(asymptotic regime), with explicit threshold-baseline bound
$\Tbeta \geq A/\deltastar$ for the non-asymptotic case.

\textbf{Round C other issues.}
\begin{itemize}
\item (C11.CLK) too weak: existence of $N_{\min}, q_{\min} > 0$
  doesn't imply $\Tbeta \leq \tau_{\mathrm{meta}}$ with calibrated
  probability.
\item $\beta'$ mislabeled: v3 used $\beta' = \exp(-\kappa
  \tau_{\mathrm{meta}} \deltastar)$ (the SPRT detection-tail form)
  for the cascade-clock event $\{\Tbeta \leq \tau_{\mathrm{meta}}\}$,
  conflating two distinct probabilities.
\item Constant taxonomy missing: $\alpha$ (Type-I rate), $R$ (LLR
  range), channel multiplicity policy, separate $\bclk$ for clock
  failure, separate from SPRT detection failure.
\item Q2 boundary convention sufficient inside Lemma 6 but must
  surface into Theorem~1's Layer~2 proof at integration time.
\end{itemize}

\textbf{Changes from v3 (Round C fixes):}

\begin{itemize}
\item \textbf{Q1 fix (algebra).}  Define $\Tbeta$ directly from
  Lemma~4's Hoeffding form:
  \[
    \Tbeta(\beta, \deltastar) \;:=\;
    \max\!\left\{\frac{2A}{\deltastar},\;
    \frac{R^2 \log(1/\beta)}{2 \deltastar^2}\right\},
  \]
  where $A := \log((1-\beta)/\alpha)$ is Paper~5's SPRT threshold
  and $R$ is the LLR-increment range from (C5.HOEFF).  In the
  asymptotic regime $\Tbeta \gg A/\deltastar$, the second term
  dominates and gives $\Pr[T_{\mathrm{detect}} > \Tbeta] \leq
  \beta$.  Critical: $\Tbeta$ scales as $1/\deltastar^2$, not
  $1/\deltastar$.
\item \textbf{(C5.HOEFF) added.}  Bounded LLR-increment range
  clause: per-step SPRT log-likelihood-ratio increments are bounded
  in absolute value by $R$, with $R$ deployment-class (not scaling
  with $|P|$ via channel multiplicity).
\item \textbf{(C11.CLK) strengthened.}  Replaced existence-only
  formulation with: there exists deterministic $\Ncasc$ such that
  $\Pr[\Nev(\tau_{\mathrm{meta}}) \geq \Ncasc] \geq 1 - \bclk$, and
  $\Ncasc \geq \Tbeta$ as an audit-time calibration constraint.
  This makes the bound $\Tbeta \leq \tau_{\mathrm{meta}}$
  operational with named failure probability $\bclk$.
\item \textbf{$\beta'$ relabeled.}  V3's use of $\beta' = \exp(-\kappa
  \tau_{\mathrm{meta}} \deltastar)$ is dropped from this lemma's
  statement.  The cascade-clock event $\{\Tbeta \leq
  \tau_{\mathrm{meta}}\}$ now has its own probability bound via
  $\bclk$ from (C11.CLK).  Total Layer~2 failure probability is
  $\beta + \bclk$.
\item \textbf{Constant taxonomy completed.}  $\alpha$ surfaced as
  monitor-design parameter; $R$ surfaced via (C5.HOEFF); $\bclk$
  surfaced via (C11.CLK); channel multiplicity flagged as
  deployment-class via (C5.HOEFF) range constancy.
\item \textbf{Q2 integration note.}  V3's $t_0 := t_0^-$ convention
  carried forward; integration must update Theorem~1's Layer~2
  proof to express the Layer~1$\to$Layer~2 transition in SPRT
  exposure clock.
\end{itemize}

\textbf{Carried forward from v3 (unchanged in v4):}
\begin{itemize}
\item Supremum-form $h_{\mathrm{detect}}(\theta) := \sup_s
  \epsgap(s)$ matching Theorem~1's Layer~2 application.
\item Discrete boundary convention $t_0 := t_0^-$.
\item $\rhogap$ rename and reinforcement that $\rhogap$ excludes
  $K_{\mathrm{Lip}}$.
\item Union-class KL floor $\deltastar = \min(\deltaadv,
  \delta_{\mathrm{adv}}^{\mathrm{env}})$.
\item (C11) framed for promotion to theorem-level condition;
  source-citation correction re Paper~9.
\end{itemize}

\section{Setup: the monitoring clock and $h_{\mathrm{detect}}$}

Theorem~\ref{thm:deployment_safety}'s Layer~2 covers adversarial
strategies $q \in \Aadv \cup \Aadv^{\mathrm{env}}$ during the
detection window
$[t_0, t_0 + T_{\mathrm{detect}}]$.

\paragraph{Boundary convention.}  We define $t_0$ as the
\emph{last SPRT exposure step} at which the deployment is still
within Layer~1's safe region; equivalently, $t_0$ is the SPRT step
immediately preceding the first crossing event ($t_0 := t_0^-$ in
the v2 notation).  $T_{\mathrm{detect}}$ is the number of SPRT
steps from $t_0$ until the SPRT log-likelihood ratio crosses the
detection threshold (Lemma~\ref{lem:4}); the crossing step itself
is counted within $T_{\mathrm{detect}}$.

Under this convention, at $s = s_0 = t_0$ the Layer~1 bound applies
exactly: $\epsgap(t_0) \leq h_{\mathrm{static}}(\theta)$.  The
single-step jump that first crosses Layer~1 occurs at $s_1$, and is
bounded by $\rhogap(\deltastar)$ via (C11).  No ``absorbed buffer''
is required.

\paragraph{Detection-window bound.}
Layer~2 applies $h_{\mathrm{detect}}$ as a point-in-time bound:
\[
  |g(\Ttruth) - g(\Pproxy)|(s) \;\leq\; \mathrm{Lip}(g) \cdot
  h_{\mathrm{detect}}(\theta)
  \quad \text{for all } s \in [t_0, t_0 + T_{\mathrm{detect}}].
\]
Therefore $h_{\mathrm{detect}}$ must bound $\epsgap(s)$ at every
$s$ in the window:
\[
  h_{\mathrm{detect}}(\theta) \;:=\; \sup_{s \in
  [t_0, t_0 + T_{\mathrm{detect}}]} \epsgap(s).
\]

\paragraph{The monitoring clock and (C11.CLK) clock comparability.}
The SPRT machinery (Paper~5) operates on a per-event clock: each
ledger entry triggers an SPRT update, and the per-step KL divergence
drives detection.  All times in this lemma — $t_0$,
$T_{\mathrm{detect}}$, $\Tbeta$, and the cascade-time references
$\tau_{\mathrm{meta}}$, $T_{\mathrm{cascade}}$ — are expressed in
this SPRT exposure clock.

\begin{assumption}[(C5.HOEFF): Bounded LLR-increment range]
\label{ass:C5HOEFF}
The per-step SPRT log-likelihood-ratio increments
$\ell_n := \log(p_1(X_n)/p_0(X_n))$ are bounded in absolute value
by a deployment-class constant $R$:
\[
  |\ell_n| \;\leq\; R \quad \text{almost surely, for all } n,
\]
with $R$ independent of $|P|$.  The deployment's channel-monitoring
policy (i.e., the number of monitored channels and the per-channel
likelihood-ratio computation) is calibrated such that $R$ does not
grow with capability count: the channel partition is either fixed
(four-channel agent-side observable plus environment-side mirror)
or family-wise corrected to keep $R$ bounded.
\end{assumption}

\begin{assumption}[(C11.CLK): Clock comparability with calibrated probability]
\label{ass:C11CLK}
There exists a deterministic deployment-class constant $\Ncasc$ and
a calibrated clock-failure probability $\bclk \in (0,1)$ such that:
\[
  \Pr\bigl[\Nev(\tau_{\mathrm{meta}}) \geq \Ncasc\bigr] \;\geq\;
  1 - \bclk,
\]
where $\Nev(\cdot)$ is the SPRT exposure event count up to wall-clock
time $\cdot$.  Furthermore, the operator calibrates $\Ncasc$ such
that the audit-time inequality
\[
  \Ncasc \;\geq\; \Tbeta(\beta, \deltastar)
\]
holds, where $\Tbeta$ is the SPRT detection quantile defined below.
\end{assumption}

\textbf{Why (C11.CLK) is strengthened in v4.}  V3's existence-only
form ($N_{\min}, q_{\min}$) did not chain to the
$\Tbeta \leq \tau_{\mathrm{meta}}$ inequality: it only said
\emph{some} number of events occur with \emph{some} probability.
V4's form combines a deterministic floor ($\Ncasc$, audit-calibrated
to exceed $\Tbeta$) with a named clock-failure probability
($\bclk$).  This makes the cascade-clock event $\{\Tbeta \leq
\tau_{\mathrm{meta}}\}$ have probability at least $1 - \bclk$
\emph{operationally}, not just by appeal to bounded event throughput.

\textbf{Why (C5.HOEFF) is required.}  Without bounded LLR-increment
range, Lemma~4's Hoeffding step
($\Pr[\Lambda_t < t\delta - s] \leq \exp(-2s^2/(tR^2))$) doesn't
apply, and the $R$-dependence of $\Tbeta$ would not be a
deployment-class constant.  In particular, channel multiplicity
scaling with $|P|$ would inflate $R$ and break intensivity.

\section{Bounded gap-growth-rate condition (C11)}

\begin{assumption}[(C11): Bounded gap-growth rate per SPRT step]
\label{ass:C11}
For any adversarial strategy $q \in \Aadv \cup
\Aadv^{\mathrm{env}}$ producing channel deviations bounded by
$\deltastar = \min(\deltaadv, \delta_{\mathrm{adv}}^{\mathrm{env}})$
(the Lemmas~5b and~5e KL floors), the per-SPRT-step proxy-truth
gap growth is bounded:
\[
  (\Delta \epsgap)_n \;:=\; \epsgap(s_n) - \epsgap(s_{n-1}) \;\leq\;
  \rhogap(\deltastar),
\]
where $\rhogap(\deltastar)$ is a deployment-class constant
depending on $\mathcal{D}$'s operational parameters (channel
sensitivity-to-gap mapping, event-rate normalization) but
independent of $|P|$.

\textbf{Continuous-time variant.}  Equivalently, in continuous
SPRT exposure clock $s$:
\[
  \frac{d \epsgap}{ds} \;\leq\; \rhogap(\deltastar).
\]
The discrete and continuous forms agree under fixed step duration
or bounded event-rate normalization.
\end{assumption}

\textbf{Operational interpretation.}  Channel deviations bounded by
$\deltastar$ (else SPRT fires faster) limit how much $\epsgap$ can
grow per SPRT step.  The constant $\rhogap$ characterizes the
deployment's channel-deviation-to-gap-growth coupling.

\textbf{$\rhogap$ is calibrated in the $\epsgap$ metric only.}
$\rhogap$ is a bound on the per-step proxy-truth gap process
$\epsgap = $ (Microfoundation's relative sup-norm), \emph{before}
applying any alignment-property mapping $g$.  In particular,
$\rhogap$ excludes the Lipschitz constant $K_{\mathrm{Lip}}$ of $g$;
the Layer~2 operational form multiplies $h_{\mathrm{detect}}$ by
$K_{\mathrm{Lip}}$ separately via (C6).  This eliminates any
double-counting between (C11) gap-dynamics and (C6) alignment-side
smoothness.

\textbf{(C11) is new content.}  \citep[Microfoundation]{lasser2026micro}'s
``$T$ failure modes'' remark notes adequacy/failure modes of the
operational truth $T$ but does not establish a formal per-step
gap-growth-rate bound.  (C11) is a Paper~10 operational assumption,
analogous to (C7) bounded co-evolution and SA1 HHI surrogate
adequacy: deployment-class condition with empirical testability,
not a derived consequence of source-paper machinery.

\section{Statement of Lemma 6 (v4)}

\begin{lemma}[$h_{\mathrm{detect}}$ intensivity]
\label{lem:6}
Under Theorem~\ref{thm:deployment_safety}'s conditions (C5)
continuous SPRT monitoring, (C5.HOEFF) bounded LLR-increment range,
(C11) bounded gap-growth rate, and (C11.CLK) clock comparability
with calibrated probability, combined with Lemma~\ref{lem:4}
(SPRT Wald--Hoeffding tail bound) and Lemma~\ref{lem:1}
(Layer~1 intensivity):

\textbf{Detection quantile.}  Define the SPRT high-probability
detection quantile derived from Lemma~4's exact Hoeffding form:
\[
  \Tbeta(\beta, \deltastar) \;:=\;
  \max\!\left\{\frac{2A}{\deltastar},\;
  \frac{R^2 \log(1/\beta)}{2 \deltastar^2}\right\},
\]
where $A := \log((1-\beta)/\alpha)$ is Paper~5's SPRT threshold
(monitor-design parameter), $R$ is the LLR-increment range
((C5.HOEFF)), $\beta$ is the target detection-tail
failure probability, and $\deltastar = \min(\deltaadv,
\delta_{\mathrm{adv}}^{\mathrm{env}})$ is the union-class KL floor
(Lemmas~\ref{lem:5b}, \ref{lem:5e}).

(i) [Tail bound on $T_{\mathrm{detect}}$.]  In the asymptotic
regime $\Tbeta \gg A/\deltastar$ (where the second term dominates),
Lemma~\ref{lem:4}'s Hoeffding form gives:
\[
  \Pr[T_{\mathrm{detect}} > \Tbeta] \;\leq\; \exp(-2 \Tbeta
  \deltastar^2 / R^2) \;=\; \beta.
\]

(ii) [Detection-window supremum.]  On the event
$\{T_{\mathrm{detect}} \leq \Tbeta\}$ (probability at least
$1-\beta$):
\[
  h_{\mathrm{detect}}(\theta) \;:=\;
  \sup_{s \in [t_0, t_0 + T_{\mathrm{detect}}]} \epsgap(s)
  \;\leq\; h_{\mathrm{static}}(\theta) + \rhogap(\deltastar) \cdot
  \Tbeta.
\]

(iii) [Lead-time-before-cascade.]  By (C11.CLK)'s deterministic
floor $\Ncasc \geq \Tbeta$ and clock-failure probability $\bclk$:
\[
  \Pr\bigl[\Tbeta \leq \Nev(\tau_{\mathrm{meta}})\bigr] \;\geq\; 1 -
  \bclk,
\]
i.e., $\Tbeta \leq \tau_{\mathrm{meta}}$ in SPRT exposure clock
with probability at least $1 - \bclk$.

(iv) [Layer~2 total failure bound.]  Total Layer~2 failure
probability (detection failure or cascade-clock miscalibration):
\[
  \Pr[\text{Layer~2 fails}] \;\leq\; \beta + \bclk.
\]

\textbf{Intensivity.}  $h_{\mathrm{detect}}(\theta)$ is intensive
in $|P|$ over the deployment class $\mathcal{D}$:
\begin{itemize}
\item $h_{\mathrm{static}}(\theta)$ is intensive by Lemma~\ref{lem:1}.
\item $\rhogap(\deltastar)$ is intensive by (C11).
\item $\Tbeta = R^2 \log(1/\beta)/(2 \deltastar^2)$ is intensive
  because $R$ is deployment-class by (C5.HOEFF), $\deltastar$ is
  deployment-class by (Lemmas~5b/5e), and $\beta$ is a monitor-design
  parameter.
\item Total bound $h_{\mathrm{static}} + \rhogap \cdot \Tbeta$ is
  a sum of intensive quantities.
\end{itemize}
\end{lemma}

\section{Proof}

\begin{proof}
We bound $h_{\mathrm{detect}}$ as the supremum of $\epsgap$ over
the detection window, derive $\Tbeta$ directly from Lemma~4's
Hoeffding form, and apply (C11.CLK) to bound the cascade-clock
event.

\emph{Step 1 (per-step decomposition).}  Under the boundary
convention from §2, set $s_0 := t_0$ (the last safe SPRT step) and
$s_n := t_0 + n$ in SPRT exposure clock.  For any
$0 \leq n \leq T_{\mathrm{detect}}$:
\[
  \epsgap(s_n) \;=\; \epsgap(s_0) + \sum_{i=1}^{n}
  (\Delta \epsgap)_i.
\]

\emph{Step 2 (apply (C11)).}  By Assumption~\ref{ass:C11}, each
per-step increment satisfies $(\Delta \epsgap)_n \leq
\rhogap(\deltastar)$.  Therefore:
\[
  \epsgap(s_n) \;\leq\; \epsgap(s_0) + n \cdot
  \rhogap(\deltastar).
\]

\emph{Step 3 (boundary condition).}  By the boundary convention,
$s_0 = t_0$ is the last SPRT step at which the deployment is still
within Layer~1's safe region.  Layer~1's intensivity result
(Lemma~\ref{lem:1}) gives:
\[
  \epsgap(s_0) \;=\; \epsgap(t_0) \;\leq\;
  h_{\mathrm{static}}(\theta).
\]
The first crossing event occurs at $s_1$, where (C11) bounds the
single-step jump by $\rhogap(\deltastar)$.

\emph{Step 4 (supremum over window).}  Combining Steps~1--3:
\[
  h_{\mathrm{detect}}(\theta) \;=\;
  \sup_{0 \leq n \leq T_{\mathrm{detect}}} \epsgap(s_n)
  \;\leq\; h_{\mathrm{static}}(\theta) +
  T_{\mathrm{detect}} \cdot \rhogap(\deltastar).
\]

\emph{Step 5 (derive $\Tbeta$ from Lemma 4 Hoeffding form).}  By
Lemma~\ref{lem:4} (\texttt{appendices/proofs.tex} §A.4), the SPRT
detection time satisfies the exact Hoeffding form:
\[
  \Pr[T_{\mathrm{detect}} > t] \;\leq\;
  \exp\!\left(-\frac{2(t \deltastar - A)^2}{t R^2}\right)
  \quad \text{for } t > A/\deltastar.
\]
We solve $\exp(-2(t\deltastar - A)^2 / (t R^2)) \leq \beta$ for
the smallest such $t$.  In the asymptotic regime
$t \deltastar \gg A$, the exponent simplifies:
$2(t\deltastar - A)^2/(tR^2) \approx 2 t \deltastar^2/R^2$, giving
the asymptotic bound
$\exp(-2 t \deltastar^2/R^2) \leq \beta$, i.e.,
$t \geq R^2 \log(1/\beta) / (2 \deltastar^2)$.

To ensure both regimes are covered (asymptotic and threshold-baseline):
\[
  \Tbeta \;:=\;
  \max\!\left\{\frac{2A}{\deltastar},\;
  \frac{R^2 \log(1/\beta)}{2 \deltastar^2}\right\}.
\]
The first term ensures $\Tbeta > A/\deltastar$ (Hoeffding regime
condition); the second term gives the tail-quantile bound.  In the
operationally relevant regime where the second term dominates,
$\Pr[T_{\mathrm{detect}} > \Tbeta] \leq \beta$.

\emph{Step 6 (detection-window supremum bound).}  On the event
$\{T_{\mathrm{detect}} \leq \Tbeta\}$ (probability at least
$1-\beta$):
\[
  h_{\mathrm{detect}}(\theta) \;\leq\;
  h_{\mathrm{static}}(\theta) + \rhogap(\deltastar) \cdot
  \Tbeta.
\]
This proves part~(ii) of the lemma.

\emph{Step 7 (cascade-clock event via (C11.CLK)).}  Separately
from the SPRT tail, (C11.CLK) provides:
$\Pr[\Nev(\tau_{\mathrm{meta}}) \geq \Ncasc] \geq 1 - \bclk$, with
audit constraint $\Ncasc \geq \Tbeta$.  Therefore:
\[
  \Pr[\Tbeta \leq \Nev(\tau_{\mathrm{meta}})] \;\geq\;
  \Pr[\Tbeta \leq \Ncasc \leq \Nev(\tau_{\mathrm{meta}})]
  \;\geq\; 1 - \bclk.
\]
This proves part~(iii) without conflating the SPRT detection-tail
form with the cascade-clock event (the v3 mistake).

\emph{Step 8 (Layer 2 total failure bound).}  By union bound on
the two failure modes (detection-tail $\{T_{\mathrm{detect}} >
\Tbeta\}$ with probability $\leq \beta$; clock-failure
$\{\Tbeta > \Nev(\tau_{\mathrm{meta}})\}$ with probability
$\leq \bclk$):
\[
  \Pr[\text{Layer~2 fails}] \;\leq\; \beta + \bclk.
\]
This proves part~(iv).

\emph{Step 9 (intensivity).}  Each constant is deployment-class
independent of $|P|$:
\begin{itemize}
\item $h_{\mathrm{static}}(\theta)$: intensive by Lemma~\ref{lem:1}
  (under (C6)--(C10)).
\item $\rhogap(\deltastar)$: deployment-class by (C11).
\item $\Tbeta = R^2 \log(1/\beta)/(2\deltastar^2)$ (asymptotic):
  $R$ is deployment-class by (C5.HOEFF) (channel-monitoring policy
  bounds LLR range without growth in $|P|$); $\deltastar$ is
  deployment-class via Lemmas~5b/5e (channel-restricted KL floor
  for adversarial classes); $\beta$ is a monitor-design parameter
  independent of $|P|$.
\item $\Ncasc, \bclk$: deployment-class by (C11.CLK) (operator
  audit-time calibration).
\end{itemize}
Therefore $h_{\mathrm{detect}}(\theta)$ is intensive in $|P|$ over
$\mathcal{D}$.  $\Box$
\end{proof}

\section{Discussion}

\paragraph{The two roles of cascade time (carried from v3).}
\begin{itemize}
\item \emph{Integration horizon for $\rhogap$:} $\Tbeta$, the SPRT
  high-probability detection quantile, gives a legitimate
  \emph{upper} bound on $T_{\mathrm{detect}}$ derived from
  Lemma~4's Hoeffding form.
\item \emph{Lead-time-before-cascade guarantee:}
  $\tau_{\mathrm{meta}} \leq T_{\mathrm{cascade}}$ (a \emph{lower}
  bound on cascade time, from
  \citep[Phase~Redundancy]{lasser2026phase}).  In v4, this is
  composed with $\Tbeta$ via (C11.CLK)'s deterministic floor
  $\Ncasc \geq \Tbeta$ and clock-failure probability $\bclk$ —
  \emph{not} via the SPRT detection-tail form (which was v3's
  mistake).
\end{itemize}

\paragraph{Why v3's $\beta'$ labeling was wrong.}
V3 labeled $\beta' = \exp(-\kappa \tau_{\mathrm{meta}} \deltastar)$
as the failure probability for the cascade-clock event $\{\Tbeta
\leq \tau_{\mathrm{meta}}\}$.  But this expression is the SPRT
\emph{detection-tail} form (probability that detection takes longer
than time $\tau_{\mathrm{meta}}$), not a cascade-clock
miscalibration probability.  V4 separates them: $\beta$ is the
detection-tail parameter; $\bclk$ is the clock-calibration failure
parameter (named in (C11.CLK)); their union bound gives total
Layer~2 failure $\leq \beta + \bclk$.

\paragraph{Why v3's $\Tbeta$ algebra was wrong.}
Lemma~4 gives $\Pr[T_{\mathrm{detect}} > t] \leq \exp(-2T\delta^2/R^2)$
with $\kappa := 2\delta/R^2$.  Therefore the exponent
$\kappa T \delta = 2T\delta^2/R^2$ scales as $\delta^2$, not $\delta$.
V3's formula $\Tbeta = \log(1/\beta)/(\kappa\delta)$ implicitly
treated $\kappa$ as $\delta$-independent, but $\kappa$ contains a
factor of $\delta$.  V4's correct form $\Tbeta = R^2 \log(1/\beta) /
(2 \deltastar^2)$ inverts the Hoeffding bound directly, making the
$1/\deltastar^2$ scaling explicit and surfacing $R$ via (C5.HOEFF).

\paragraph{Lemma 5d integration item (Round B Q3 follow-up,
carried).}
At integration time, §4.10's Lemma~5d should be restated using the
union-class floor $\deltastar$.  This is a small amendment to §4.10
and §A.4 that should accompany Lemma~6 integration.

\paragraph{Constant taxonomy (v4 complete).}
\begin{itemize}
\item $h_{\mathrm{static}}(\theta)$: intensive composition
  (Lemma~\ref{lem:1}).
\item $\rhogap(\deltastar)$: \emph{Paper 10 operational constant}
  via (C11); calibrated by per-channel sensitivity-to-gap
  measurement (Audit~7).
\item $\Tbeta = R^2 \log(1/\beta)/(2\deltastar^2)$ (asymptotic):
  SPRT-derived detection quantile.
\item $A = \log((1-\beta)/\alpha)$: Paper~5 SPRT threshold.
\item $\alpha$: Type-I error rate, monitor-design parameter
  (deployment-class).
\item $\beta$: Type-II error rate / detection-tail parameter,
  monitor-design parameter (deployment-class).
\item $R$: per-step LLR-increment range, surfaced via
  (C5.HOEFF) (deployment-class; channel-multiplicity policy bounds
  $R$ without growth in $|P|$).
\item $\Ncasc$: deterministic event-count floor, surfaced via
  (C11.CLK) (audit-calibrated to satisfy $\Ncasc \geq \Tbeta$).
\item $\bclk$: clock-calibration failure probability, surfaced via
  (C11.CLK).
\item $\tau_{\mathrm{meta}}, T_{\mathrm{cascade}}$: source-derived
  cascade-time references from
  \citep[Phase~Redundancy]{lasser2026phase} (lower bound only;
  with the (C7) bounded co-evolution caveat).
\end{itemize}

\paragraph{Layer 2 operational form.}
Combining Lemma~6 with $\mathrm{Lip}(g) \leq K_{\mathrm{Lip}}$ from
(C6):
\[
  |g(\Ttruth) - g(\Pproxy)|(s) \;\leq\; K_{\mathrm{Lip}} \cdot
  \bigl(h_{\mathrm{static}}(\theta) + \rhogap(\deltastar) \cdot
  \Tbeta\bigr)
\]
with probability at least $1 - (\beta + \bclk)$ over the joint
event $\{T_{\mathrm{detect}} \leq \Tbeta\} \cap \{\Tbeta \leq
\Nev(\tau_{\mathrm{meta}})\}$.  The right-hand side is the sum of
intensive constants times $K_{\mathrm{Lip}}$, all $|P|$-independent.

\section{What this proof does and does not establish}

\textbf{Establishes:}
\begin{itemize}
\item Algebraically correct $\Tbeta$ derivation from Lemma~4's
  Hoeffding form, with explicit $1/\deltastar^2$ scaling.
\item (C5.HOEFF) bounded LLR-increment range as a deployment-class
  condition, ensuring $R$ does not scale with $|P|$ via channel
  multiplicity.
\item (C11.CLK) strengthened with deterministic floor $\Ncasc \geq
  \Tbeta$ and named clock-failure probability $\bclk$, separating
  cascade-clock event from SPRT detection-tail event.
\item Supremum-form $h_{\mathrm{detect}}$ definition matching
  Theorem 1's Layer 2 application.
\item Rigorous discrete-time boundary condition under the
  $t_0 := t_0^-$ convention.
\item Union-class KL floor $\deltastar$ for Layer 2's full
  detection class.
\item Intensive bound on $h_{\mathrm{detect}}$ as $h_{\mathrm{static}}
  + \rhogap \Tbeta$ with high probability $1-\beta$.
\item Layer~2 total failure probability $\leq \beta + \bclk$ via
  union bound.
\item Constant taxonomy with all $|P|$-dependent constants surfaced
  ($\alpha, \beta, R, \Ncasc, \bclk$ all deployment-class via
  named conditions).
\end{itemize}

\textbf{Does not establish (deferred):}
\begin{itemize}
\item Tightness of $\rhogap(\deltastar), R, \Ncasc$.  Operators
  may calibrate tighter bounds for restricted threat models.
\item Existence of deployment configurations satisfying (C5.HOEFF),
  (C11), and (C11.CLK).  Verified empirically (Audit~7).
\item Structural derivation of (C11) from finer source-paper
  machinery.  Currently a Paper 10 operational assumption.
\item Layer~2 proof updates in Theorem~1 (must adopt $t_0 := t_0^-$
  convention and use $\Tbeta$ instead of $T_{\mathrm{cascade}}$ in
  Step~6 of the Layer~2 proof).  This is an integration task.
\end{itemize}

\section{Specific verification questions for v4 codex review}

\begin{enumerate}
\item \textbf{(Q1 algebra fix.)}  Is the v4 derivation of $\Tbeta =
  R^2 \log(1/\beta)/(2\deltastar^2)$ from Lemma~4's Hoeffding form
  correct?  Step~5 of the proof inverts the asymptotic Hoeffding
  bound $\exp(-2T\deltastar^2/R^2) \leq \beta$.  The
  $\max\{2A/\deltastar, R^2\log(1/\beta)/(2\deltastar^2)\}$ form
  ensures both the asymptotic-regime condition
  ($t > A/\deltastar$) and the tail-quantile bound.  Is this
  derivation rigorous, and is the $1/\deltastar^2$ scaling
  correctly captured?  In particular, is the asymptotic
  approximation $2(t\deltastar - A)^2/(tR^2) \approx 2t\deltastar^2/R^2$
  sufficient or does the lemma need the exact form?

\item \textbf{((C5.HOEFF) sufficiency.)}  Is (C5.HOEFF) bounded
  LLR-increment range the right framing to make $R$ deployment-class?
  Does it correctly handle: (a) fixed four-channel monitoring (Paper
  5's default), (b) environment-side mirror channels (Paper 10's
  Lemma 5e extension), (c) potential multinomial channel partitions?
  Is the family-wise correction clause tight enough?

\item \textbf{((C11.CLK) strengthening.)}  Is the v4 form
  ($\Pr[\Nev(\tau_{\mathrm{meta}}) \geq \Ncasc] \geq 1 - \bclk$ +
  audit constraint $\Ncasc \geq \Tbeta$) sufficient to derive
  $\Pr[\Tbeta \leq \tau_{\mathrm{meta}}] \geq 1 - \bclk$ rigorously?
  Step~7 uses the chain $\Tbeta \leq \Ncasc \leq
  \Nev(\tau_{\mathrm{meta}})$; is this watertight, or does the
  audit-time inequality $\Ncasc \geq \Tbeta$ need probabilistic
  treatment as well (e.g., if calibration error propagates)?

\item \textbf{(Probability-labeling.)}  V4 uses $\beta$ for
  detection-tail and $\bclk$ for clock-failure, with total Layer~2
  failure $\leq \beta + \bclk$ via union bound.  Is this labeling
  now rigorous and unambiguous, or are there remaining places where
  $\beta'$ from v3 (the discarded SPRT-detection-tail variant for
  cascade-clock) might re-enter through Lemma~5d composition?

\item \textbf{(Constant taxonomy completeness.)}  V4's Discussion
  enumerates $h_{\mathrm{static}}, \rhogap, \Tbeta, A, \alpha,
  \beta, R, \Ncasc, \bclk, \tau_{\mathrm{meta}}, T_{\mathrm{cascade}}$.
  Are there any remaining hidden $|P|$-dependent constants?  In
  particular: is $\alpha$'s deployment-class status sufficient (it
  affects $A = \log((1-\beta)/\alpha)$), or does $\alpha$ also
  need a named clause?  Is the channel-multiplicity policy
  adequately surfaced via (C5.HOEFF), or does it need a separate
  (C5.MULT) clause?
\end{enumerate}

\end{document}
